\documentclass[11pt]{amsart}
\usepackage[margin=1.15in]{geometry}
\usepackage{amsmath,amssymb,amsthm}
\usepackage[hidelinks]{hyperref}

\newtheorem{theorem}{Theorem}
\newtheorem{proposition}[theorem]{Proposition}
\newtheorem{lemma}[theorem]{Lemma}
\theoremstyle{remark}
\newtheorem{remark}[theorem]{Remark}

\newcommand{\Z}{\mathbb{Z}}
\newcommand{\N}{\mathbb{N}}
\newcommand{\e}{\mathrm{e}}
\newcommand{\lean}[1]{\texttt{#1}}
\DeclareMathOperator{\re}{Re}

\title[Prime sums, prime counting and prime finding in time $N^{o(1)}$]%
{Exact prime sums, exact prime counting and deterministic prime finding
in time $N^{o(1)}$, relative to the new counting engines}

\author{Jiayan Ni}
\address{Cornell University, Ithaca, NY 14853, United States}
\email{jn577@cornell.edu}
\date{November 15, 2026. Announcement; a full paper is in preparation.}

\begin{document}

\begin{abstract}
We announce three machine-checked theorems.  There are deterministic Turing
machines that, on input $N$ in binary, output (i)~the exact sum
$S(N)=\sum_{p\le N}p$ of the primes up to $N$, (ii)~the exact prime count
$\pi(N)$, and (iii)~for $N\ge1$ a prime $p$ with $N<p\le 2N$, each within
$2^{O(b/\sqrt{\log b})}=N^{o(1)}$ steps, where $b$ is the bit length of $N$.
The best previous bounds for all three problems were of the form
$N^{1/2+o(1)}$.  The proofs are reductions to the subexponential counting
algorithm behind the recent refutation of the Exponential Time Hypothesis:
they are not classical algorithms, they are purely asymptotic, and they stand
or fall with that result.  The new ingredient is an analytic one: Ramanujan
sums express the M\"obius function without factoring, and a Cauchy-kernel
quadrature turns the phases, the logarithm and its reciprocal into short
linear combinations of rational functions that a counting network of linear
size can evaluate.  All three statements are formalized in Lean~4 against a
specification file that imports only Mathlib.
\end{abstract}

\maketitle

\section{Results}

For a natural number $N$ let $b(N)$ be the length of its binary expansion
(so $b(0)=0$), and for a natural number $C$ put
\[
  T_C(N)=\Bigl\lfloor C\cdot 2^{\,C\,b(N)/\sqrt{\log_2(b(N)+2)}}\Bigr\rfloor .
\]
Numbers are written in binary, least significant bit first, and $0$ is the
empty word.  ``Machine'' means a deterministic multi-stack Turing machine
with finite stack alphabets, in the sense of Mathlib's \lean{TM2} model with
its timed halting relation \lean{TM2OutputsInTime}; the step count includes
everything the machine does.

\begin{theorem}[Exact prime sum]\label{thm:sum}
There are a machine $M$ and a constant $C$ such that for every $N\in\N$, on
input $N$ the machine $M$ halts within $T_C(N)$ steps with output
$S(N)=\sum_{p\le N,\ p\ \mathrm{prime}}p$.
\end{theorem}

\begin{theorem}[Exact prime counting]\label{thm:count}
There are a machine $M$ and a constant $C$ such that for every $N\in\N$, on
input $N$ the machine $M$ halts within $T_C(N)$ steps with output $\pi(N)$.
\end{theorem}

\begin{theorem}[Deterministic prime finding]\label{thm:find}
There are a machine $M$ and a constant $C$ such that for every $N\in\N$, on
input $N$ the machine $M$ halts within $T_C(N)$ steps and outputs a prime $p$
with $N<p\le 2\max(N,1)$.
\end{theorem}

In each theorem one machine and one constant serve all inputs, including
$N=0$ and $N=1$; the three theorems may use different machines and constants.
Since $b(N)=\lfloor\log_2 N\rfloor+1$ for $N\ge1$, the bound is
$2^{O(\log N/\sqrt{\log\log N})}=N^{o(1)}$.

\subsection*{Context}
The analytic method of Lagarias and Odlyzko \cite{LO} computes $\pi(N)$ in
time $N^{1/2+o(1)}$, and the elementary method of Hirsch, Kessler and
Mendlovic \cite{HKM} gives $\widetilde O(\sqrt N)$ for both $\pi(N)$ and
$S(N)$; the combinatorial methods \cite{LMO,DR} run in time about $N^{2/3}$.
For deterministically finding a prime larger than $N$ the best known bound is
also $N^{1/2+o(1)}$, and improving it was the goal of the Polymath4 project
\cite{TCH}.  We know of no earlier bound of the form $N^{o(1)}$ for any of
the three problems, conditional or not.

In October 2026 two deterministic algorithms for \#3-SAT in time $2^{o(n)}$
were announced, with Lean formalizations \cite{O,A}; we use the one of
\cite{A}.
Several consequences for number-theoretic problems followed quickly, among
them deterministic factoring in time $N^{o(1)}$ \cite{RR} and, closest to the
present note, Tulsyan's computation of the \emph{parity} of $\pi(N)$ in time
$N^{o(1)}$ \cite{T1}.  The notes accompanying \cite{T1} record exact
$\pi(N)$ and deterministic prime finding as out of reach of that approach
without a further, analytic component (see also \cite{T2}).
Theorems~\ref{thm:sum}--\ref{thm:find} supply such a component.  We searched
the Hexagon server, arXiv and the public repositories of \cite{A,T1} on
November~15, 2026 and found no overlapping claim.

\subsection*{What kind of result this is}
The theorems are consequences of the counting algorithm of \cite{A}.  They
contain no unproved hypothesis in the formal sense: the Lean development of
\cite{A} is imported and checked together with ours.  But the speed comes
entirely from that algorithm, and if it were withdrawn the results would go
with it.  The constant $C$ is not extracted.  The constants that arise in
the proof are very large (the reduction alone introduces a network constant
$2^{22}\cdot132$ and a padding factor $2049$), and we make no claim about
any feasible range of $N$.  Nothing here is a practical algorithm.

\section{The engine}

We use the following consequence of \cite{A}, in the packaged form developed
in \cite{T1} and extended here to batches of queries.

A \emph{counting package} is a uniform family of systems of constraints over
bounded natural-number variables.  For an input word $w$ of length $n$ and a
query word $\pi$ of length at most $n$, the system has $O(n)$ shared bits in
total and is built from additions, comparisons, copies and a bounded number
of multiplications; it is produced from $(w,\pi)$ in polynomial time.  The
quantity of interest is its number of solutions $Q(w,\pi)$.

\begin{proposition}[Batch counting]\label{prop:engine}
Let a counting package be given, together with polynomial-time procedures
that produce polynomially many queries $\pi_1,\pi_2,\dots$ (possibly
adaptively, each depending on earlier answers) and compute an output from
$w$ and the answers $Q(w,\pi_i)$.  Then one machine computes that output in
time $2^{O(n/\sqrt{\log n})}$.  The same holds when $w$ is first replaced by
a padded word of length $O(n)$ that is computable in polynomial time.
\end{proposition}

The essential point is that the exponent is governed by the number of shared
bits of a \emph{single} system, which must be linear in $n$, while the number
of queries and all pre- and post-processing may be polynomial.
In Lean this is \lean{NatBatchQuery.machine\_TM\_linear\_input}, proved from
\lean{CountedRoute.onTheCountedRecursion} of \cite{A}.

\section{Sketch of the reduction}

\subsection{A formula without factoring}
Write $\e(t)=\exp(2\pi i t)$.  For $n\ge2$,
\[
  \mathbf 1_{\mathrm{prime}}(n)=-\frac{\mu(n)\Lambda(n)}{\log n},
\]
since $\mu(n)\Lambda(n)$ is $-\log n$ at primes and vanishes elsewhere.
With $\Lambda(n)=\sum_{dm=n}\mu(d)\log m$ and Ramanujan's evaluation
\cite{R} of $c_d(1)$,
\[
  \mu(d)=\sum_{\substack{0\le x<d\\ \gcd(x,d)=1}}\e(x/d),
\]
we obtain
\begin{equation}\label{eq:four}
  S(N)=-\sum_{\substack{dm=n\\ 2\le n\le N}}\;
  \sum_{\substack{0\le x<d,\ 0\le x'<n\\ (x,d)=(x',n)=1}}
  n\cdot\e\!\Bigl(\frac{x'}{n}\Bigr)\cdot\frac{1}{\log n}\cdot
  \e\!\Bigl(\frac{x}{d}\Bigr)\cdot\log m .
\end{equation}
Each summand is $n$ times a product of four values of analytic functions at
rational points with denominators $n$, $d$ or a power of two.  No factor
requires knowing the factorization of $n$.  Coprimality is expressed by a
unique bounded B\'ezout witness: for $d\ge2$ and $1\le x<d$ one has
$\gcd(x,d)=1$ if and only if there is exactly one pair $(u,v)$ with
$0\le u<d$, $0\le v<x$ and $xu=1+dv$.

\subsection{Cauchy kernels}
Let $c=\tfrac32$, $r=\tfrac34$ and $z_\ell=c+r\,\e(\ell/M)$ for
$0\le\ell<M$.  For $h$ analytic on the closed disc $|z-c|\le\tfrac78$ put
\[
  Q_Mh(y)=\frac1M\sum_{\ell=0}^{M-1}h(z_\ell)\,\frac{z_\ell-c}{z_\ell-y}.
\]

\begin{lemma}\label{lem:quad}
If $|h|\le H$ on the disc $|z-c|\le\tfrac78$ and $M\ge1$, then
$|Q_Mh(y)-h(y)|\le 49\,H\,(6/7)^M$ for all real $1\le y\le2$.
\end{lemma}

This is the geometric convergence of the trapezoidal rule on a circle
\cite{TW}, with explicit constants.  We apply it to three functions: the
phase $h(z)=\e(z-1)$ at $y=1+x/d$; the logarithm
$h_j(z)=j\log2+\log z$ at $y=t/2^{j}$, where $t=2^jy$ and
$j=\lfloor\log_2t\rfloor$; and its reciprocal $1/h_j$ for $j\ge1$.  All
three are analytic and suitably bounded on the disc.  Consequently every
factor in \eqref{eq:four} is, up to an error $2^{-\Omega(M)}$, a linear
combination of $M$ \emph{reciprocal kernels} $1/(z_\ell-y)$ with coefficients
that depend on $\ell$ and on the dyadic scale $j$ only.  Using the kernel for
the phases as well avoids any long chain of rotations.

\subsection{Fixed point arithmetic and exact rounding}
Let $\beta=\lfloor\log_2N\rfloor+1$ (and $\beta=1$ for $N=0$).  We take
$M=35\beta+185$ nodes, round nodes and coefficients to $p=7\beta+32$ binary
places, and truncate every reciprocal kernel towards zero at $t=7\beta+25$
binary places.  If $z=(U+iV)/2^p$ and $y=a/q$, then
with $X=Uq-2^pa$, $Y=Vq$ and $D=X^2+Y^2>0$ the absolute values of the real
and imaginary parts of the truncated kernel are the integers
\begin{equation}\label{eq:quot}
  \Bigl\lfloor\frac{2^{t}2^{p}q\,|X|}{D}\Bigr\rfloor,\qquad
  \Bigl\lfloor\frac{2^{t}2^{p}q\,|Y|}{D}\Bigr\rfloor .
\end{equation}
Each of the four factors is then within $\delta=2^{-(6\beta+6)}$ of its true
value, a product of four is within $(12+10\beta)\delta$, and the weights and
the number of summands contribute at most $N^4$.  The total error of the
resulting finite expression $\widetilde S(N)$ is therefore at most
$N^4(12+10\beta)\,2^{-(6\beta+6)}<\tfrac12$, so that
\[
  \bigl\lfloor\re\widetilde S(N)+\tfrac12\bigr\rfloor=S(N)
  \qquad\text{for every }N\in\N .
\]
The coefficients need not be exact roundings of transcendental numbers: any
table within a stated tolerance works, and such a table is produced in
polynomial time by rational Taylor series with explicit error certificates.

\subsection{One counting network}
Expanding the four linear combinations, the real and imaginary parts and the
signs, $\widetilde S(N)$ becomes an integer linear combination, with
$\beta^{O(1)}$ terms and coefficients computed outside the engine, of sums
\[
  \sum n\cdot q_1q_2q_3q_4
\]
over the arithmetic data $(n,d,m,x,x')$ in fixed dyadic ranges, where the
$q_i$ are quotients of the form \eqref{eq:quot} at four chosen nodes.  Each
such sum is the number of solutions of one explicit system: the relation
$n=dm$, the two B\'ezout witnesses, a quotient and remainder for each kernel,
and free variables $v_i<q_i$ and $v_0<n$ that turn weights into
multiplicities.  The system has $72$ variables, $33$ equations and $132$
product slots, each variable has $O(\beta)$ bits, and all auxiliary variables are
uniquely determined.  The boundary cases ($d=1$, $x=1$, zero quotients, sign
branches) are treated separately.

\subsection{The machines}
The machine for Theorem~\ref{thm:sum} generates the coefficient table,
enumerates the queries, obtains the counts from
Proposition~\ref{prop:engine}, forms the integer linear combination, divides
exactly by the common denominator and rounds.  For Theorem~\ref{thm:count}
the multiplicity variable $v_0<n$ is replaced by $v_0<1$; the error budget
and the rest of the construction are unchanged.  For
Theorem~\ref{thm:find}, note that for $a\le c$ one has $S(a)<S(c)$ if and
only if the interval $(a,c]$ contains a prime.  By Bertrand's postulate
$(N,2\max(N,1)]$ contains a prime, and bisection with the invariant
$S(\mathrm{left})=S(N)<S(\mathrm{right})$ ends at a prime after $b(N)+3$
evaluations of $S$.  The adaptive version of
Proposition~\ref{prop:engine} accounts for all of them in one machine.

\section{Formal verification}

The three theorems are stated in a file \lean{PrimeIndependent/MainResults.lean}
that imports only Mathlib and uses none of the definitions of the proof; in
particular $\pi(N)$ is Mathlib's \lean{Nat.primeCounting} and the prime found
in Theorem~\ref{thm:find} is specified by its properties, not by the search
procedure.  The correspondence is as follows.

\smallskip
\begin{center}
\begin{tabular}{ll}
Theorem~\ref{thm:sum} & \lean{PrimeIndependent.primeSum\_time} \\
Theorem~\ref{thm:count} & \lean{PrimeIndependent.primeCounting\_time} \\
Theorem~\ref{thm:find} & \lean{PrimeIndependent.primeFinding\_time} \\
Lemma~\ref{lem:quad} & \lean{PrimeAlgorithm.Cauchy.cauchyApprox\_error} \\
Exact rounding & \lean{FixedCauchyPrimeSum.fixedPrimeSum\_all\_inputs} \\
Proposition~\ref{prop:engine} &
  \lean{NatBatchQuery.machine\_TM\_linear\_input} \\
Coefficient generation & \lean{CertifiedGenerator.generate\_certified} \\
\end{tabular}
\end{center}
\smallskip

The development uses Lean~4.33.1 and builds on the formalization of
\cite{A} (commit \lean{16a5ce4}) and on the library of \cite{T1} (commit
\lean{4193596}), neither of which is modified.  The three theorems depend
only on the axioms \lean{propext}, \lean{Classical.choice} and
\lean{Quot.sound}.  The whole dependency closure, including all of Mathlib,
was built from source without caches; the result was replayed with
\lean{leanchecker --fresh} and compared against the specification with
Comparator, which checks the proofs with the Lean kernel and with the
independent kernel nanoda.  These checks were all run on one machine;
independent reproduction is welcome.

\smallskip\noindent
Source: \url{https://github.com/6teamFurina/prime-algorithm},
release \texttt{v0.1-verified}.  The release tag identifies the verified
proof-source commit.  The repository is publicly accessible.

\section{Remarks}

\begin{remark}
The exponent in Proposition~\ref{prop:engine} depends on the number of shared
bits of one system, so the reduction must keep that number linear in $b$.
This is why the logarithm is not computed by a circuit inside the network:
it is replaced by $O(b)$ rational kernels of constant arithmetic complexity,
whose linear combination is formed outside.  The same reason forces the
input to the engine to be padded only linearly.
\end{remark}

\begin{remark}
The method computes weighted sums over the integers up to $N$ of products of
a bounded number of analytic factors and Ramanujan phases.  We expect it to
apply to other summatory functions, such as the Mertens function, but only
the three theorems above have been formalized.
\end{remark}

\begin{remark}
The full paper will contain complete proofs of Lemma~\ref{lem:quad}, of the
error budget, of the counting bijection for the network and of the bit
bounds, following the formal development.
\end{remark}

\subsection*{Release preparation}
AI tools assisted with preparing the repository and its reproduction scripts.
The mathematical verification described above is based on the recorded
checker executions.

\begin{thebibliography}{99}

\bibitem{A} Anthropic,
\emph{A $2^{o(n)}$ time algorithm for 3-SAT},
Hexagon 2610.00577 (2026); Lean formalization at
\url{https://github.com/anthropics/formal-math}.

\bibitem{DR} M.~Del\'eglise and J.~Rivat,
\emph{Computing $\pi(x)$: the Meissel, Lehmer, Lagarias, Miller, Odlyzko
method}, Math. Comp. \textbf{65} (1996), 235--245.

\bibitem{HKM} D.~Hirsch, I.~Kessler and U.~Mendlovic,
\emph{Computing $\pi(N)$: an elementary approach in $\tilde O(\sqrt N)$
time}, arXiv:2212.09857.

\bibitem{LMO} J.~C.~Lagarias, V.~S.~Miller and A.~M.~Odlyzko,
\emph{Computing $\pi(x)$: the Meissel--Lehmer method},
Math. Comp. \textbf{44} (1985), 537--560.

\bibitem{LO} J.~C.~Lagarias and A.~M.~Odlyzko,
\emph{Computing $\pi(x)$: an analytic method},
J. Algorithms \textbf{8} (1987), 173--191.

\bibitem{O} OpenAI,
\emph{The Exponential Time Hypothesis is false}, preprint (October 2026),
\url{https://github.com/openai/math}.

\bibitem{R} S.~Ramanujan,
\emph{On certain trigonometrical sums and their applications in the theory
of numbers}, Trans. Cambridge Philos. Soc. \textbf{22} (1918), 259--276.

\bibitem{RR} I.~Rudich and L.-M.~Rousseau,
\emph{Deterministic factoring in time $N^{o(1)}$},
Hexagon 2610.00600 (2026).

\bibitem{TCH} T.~Tao, E.~Croot III and H.~Helfgott,
\emph{Deterministic methods to find primes},
Math. Comp. \textbf{81} (2012), 1233--1246.

\bibitem{TW} L.~N.~Trefethen and J.~A.~C.~Weideman,
\emph{The exponentially convergent trapezoidal rule},
SIAM Rev. \textbf{56} (2014), 385--458.

\bibitem{T1} V.~Tulsyan,
\emph{The parity of the prime-counting function in time $x^{o(1)}$},
Hexagon 2611.00116 (2026); Lean library at
\url{https://github.com/xennygrimmato/eth-hamiltonicity}.

\bibitem{T2} V.~Tulsyan,
\emph{What the ETH refutations do not give: barriers for the counting
engines, with machine-checked lemmas}, Hexagon 2611.00149 (2026).

\end{thebibliography}

\end{document}
